\RequirePackage{iftex}
\ifPDFTeX
\documentclass[a4paper,12pt]{article}
\usepackage[margin=22mm]{geometry}
\usepackage[T1]{fontenc}
\usepackage[utf8]{inputenc}
\usepackage{graphicx}
\usepackage{CJKutf8}
\else
\documentclass[a4paper,12pt]{jsarticle}
\usepackage[dvipdfmx,margin=22mm]{geometry}
\usepackage[dvipdfmx]{graphicx}
\fi
\usepackage{amsmath,amssymb,amsthm}
\usepackage{enumitem}

\setlength{\parindent}{0pt}
\setlength{\parskip}{0.35em}
\setlength{\fboxsep}{8pt}
\setlist[itemize]{leftmargin=2em,itemsep=0.25em,topsep=0.35em}
\setlist[enumerate]{leftmargin=2.2em,itemsep=0.45em,topsep=0.35em}

\newcommand{\pointbox}[1]{%
\begin{center}
\fbox{\begin{minipage}{0.91\linewidth}#1\end{minipage}}
\end{center}
}

\begin{document}
\ifPDFTeX
\begin{CJK}{UTF8}{min}
\fi

\theoremstyle{definition}
\newtheorem*{definition}{定義}
\newtheorem*{theorem}{定理}
\newtheorem*{example}{例}
\newtheorem*{remark}{注意}
\renewcommand{\proofname}{\normalfont 証明}

\begin{center}
{\LARGE 数学1A 第12回レジュメ}\\[0.4em]
連続関数の基本性質
\end{center}

\section*{1. 連続性}

\begin{definition}
$D\subset\mathbb{R}$ とし、$f:D\to\mathbb{R}$ とする。$f$ が点 $a\in D$ で
連続であるとは、任意の $\varepsilon>0$ に対してある $\delta>0$ が存在し、
$x\in D$ かつ $|x-a|<\delta$ ならば
\[
|f(x)-f(a)|<\varepsilon
\]
となることをいう。
\end{definition}

\begin{theorem}[数列による連続性の判定]
$f$ が $a$ で連続であることと、$D$ の任意の数列 $x_n$ について
\[
x_n\to a\quad\Rightarrow\quad f(x_n)\to f(a)
\]
が成り立つことは同値である。
\end{theorem}

\begin{proof}
$f$ が $a$ で連続であるとする。$\varepsilon>0$ を取る。連続性より、
$|x-a|<\delta$ ならば $|f(x)-f(a)|<\varepsilon$ となる $\delta>0$ がある。
$x_n\to a$ なので、十分大きい $n$ で $|x_n-a|<\delta$ となる。したがって
$f(x_n)\to f(a)$ である。

逆に、$f$ が $a$ で連続でないとする。するとある $\varepsilon_0>0$ が存在し、
どんな $\delta>0$ に対しても $|x-a|<\delta$ かつ
$|f(x)-f(a)|\ge\varepsilon_0$ となる $x\in D$ が取れる。特に
$\delta=1/n$ として $x_n$ を取れば、$x_n\to a$ だが
$f(x_n)\to f(a)$ ではない。
\end{proof}

\section*{2. 連続関数の演算}

\begin{theorem}
$f,g$ が点 $a$ で連続ならば、$f+g$ と $fg$ も点 $a$ で連続である。
また $g(a)\ne0$ ならば $f/g$ も点 $a$ で連続である。
\end{theorem}

\begin{proof}
$x_n\to a$ とする。$f,g$ の連続性より $f(x_n)\to f(a)$,
$g(x_n)\to g(a)$ である。数列の極限の和・積・商の性質を用いれば、
\[
(f+g)(x_n)\to (f+g)(a),\qquad (fg)(x_n)\to (fg)(a)
\]
が従う。$g(a)\ne0$ のときは十分大きい $n$ で $g(x_n)\ne0$ となり、
\[
\frac{f(x_n)}{g(x_n)}\to \frac{f(a)}{g(a)}
\]
である。数列による判定から結論を得る。
\end{proof}

\begin{example}
多項式関数はすべての点で連続である。有理関数は、分母が $0$ でない点で連続である。
\end{example}

\section*{3. 中間値の定理}

\begin{theorem}[中間値の定理]
$f:[a,b]\to\mathbb{R}$ が連続で、$f(a)<\lambda<f(b)$ とする。このとき
ある $c\in(a,b)$ が存在して
\[
f(c)=\lambda
\]
となる。$f(b)<\lambda<f(a)$ の場合も同様である。
\end{theorem}

\begin{proof}
$S=\{x\in[a,b]\mid f(x)\le\lambda\}$ とおく。$a\in S$ なので $S$ は空でなく、
$b$ によって上に有界である。よって $c=\sup S$ が存在する。

連続性を用いて $f(c)=\lambda$ を示す。もし $f(c)<\lambda$ ならば、$c$ の十分近く
では $f(x)<\lambda$ である。すると $c$ より大きい $S$ の元が存在し、$c$ が上界で
あることに矛盾する。もし $f(c)>\lambda$ ならば、$c$ の十分近くでは
$f(x)>\lambda$ である。一方、上限の近似性により $c$ に左から近い $S$ の元が
存在するので矛盾する。したがって $f(c)=\lambda$ である。
\end{proof}

\section*{4. 最大値・最小値}

\begin{theorem}[最大値・最小値の定理]
$f:[a,b]\to\mathbb{R}$ が連続ならば、$f$ は $[a,b]$ 上で最大値と最小値を取る。
\end{theorem}

\begin{proof}
ここでは事実として、閉区間 $[a,b]$ の任意の数列は収束部分列をもち、その極限も
$[a,b]$ に属することを用いる。

$f$ が上に有界であるとし、$M=\sup f([a,b])$ とおく。上限の近似性より、
$f(x_n)>M-1/n$ となる $x_n\in[a,b]$ を取れる。$[a,b]$ の性質より、部分列
$x_{n_k}$ がある $x_0\in[a,b]$ に収束する。連続性より
$f(x_{n_k})\to f(x_0)$ であり、一方で $f(x_{n_k})\to M$ なので $f(x_0)=M$ である。
最小値も同様に示せる。
\end{proof}

\pointbox{
閉区間上の連続関数には、中間値を取る性質と、最大値・最小値を取る性質がある。
どちらも閉区間と連続性が重要である。
}

\ifPDFTeX
\end{CJK}
\fi
\end{document}
